\documentclass[../readings.tex]{subfiles}

\begin{document}

\subsection{Spectral Graph Theory and PageRank}
\label{sub:spectral_graph_theory}

\textBox{%
Graph algorithms become numerical linear algebra after choosing matrix conventions. Undirected graphs lead to symmetric Laplacians. Directed ranking problems lead to stochastic matrices and eigenvectors.
}

\subsubsection{Graph Matrices}

\addRemark{%
\textbf{(Convention)}\enspace For undirected graphs, $\bA$ is symmetric. For directed random walks in this section, probability vectors are columns and $P_{ij}$ is the probability of moving from node $j$ to node $i$.
}

\addDef{Adjacency and Degree Matrices}{%
For a graph $G = (V, E)$ with $n$ vertices:
\begin{itemize}
    \item \textbf{Adjacency $\bA$:}\enspace $A_{ij} = w_{ij}$ if $(i,j) \in E$, else $0$.
    \item \textbf{Degree $\bD$:}\enspace Diagonal matrix with $D_{ii} = d_i = \sum_{j} w_{ij}$.
\end{itemize}
}
\label{def:graph-adjacency-degree}

\addDef{Graph Laplacian}{%
Defined as \textbf{$\bL = \bD - \bA$}.
\begin{itemize}
    \item \textbf{Discrete Derivative:}\enspace $(\bL \boldsymbol{f})_i = \sum_{j \sim i} w_{ij} (f_i - f_j)$.
    \item \textbf{Quadratic Form:}\enspace $\bx^T \bL \bx = \sum_{(i,j) \in E} w_{ij} (x_i - x_j)^2$, with each undirected edge counted once.
\end{itemize}
}
\label{def:graph-laplacian}

\addProp{Laplacian Positive Semidefinite}{%
For an undirected weighted graph with nonnegative weights, $\bL$ is symmetric positive semidefinite and $\bL\mathbf{1}=\bzero$.
}
\label{prop:graph-laplacian-psd}

\addPrf{%
Symmetry follows from $w_{ij}=w_{ji}$. Using the quadratic form in Def~\ref{def:graph-laplacian},
\begin{align*}
    \bx^T\bL\bx
    =
    \sum_{(i,j)\in E} w_{ij}(x_i-x_j)^2
    \geq 0.
\end{align*}
Thus $\bL$ is positive semidefinite. If $\bx=\mathbf{1}$, every difference $x_i-x_j$ is zero, and direct substitution gives $\bL\mathbf{1}=\bzero$.
}

\addExmpl{%
\textbf{(Three-node path)}\enspace For $1-2-3$,
\begin{align*}
    \bA=
    \begin{pmatrix}
    0&1&0\\
    1&0&1\\
    0&1&0
    \end{pmatrix},
    \qquad
    \bL=
    \begin{pmatrix}
    1&-1&0\\
    -1&2&-1\\
    0&-1&1
    \end{pmatrix}.
\end{align*}
The zero eigenvector is $\mathbf{1}$ because constant signals have no graph variation.
}

\addThm{Spectral Properties}{%
For a connected graph with $0 = \lambda_1 \leq \lambda_2 \leq \dots \leq \lambda_n$:
\begin{enumerate}
    \item $\lambda_1 = 0$ with eigenvector $\mathbf{1}$.
    \item \textbf{Multiplicity of 0:}\enspace Equals the number of connected components.
    \item \textbf{Fiedler Value ($\lambda_2$):}\enspace Measures algebraic connectivity. $\lambda_2 > 0$ iff graph is connected.
\end{enumerate}
}
\label{thm:graph-spectral-properties}

\addExcr{%
\begin{enumerate}
    \item Use Def~\ref{def:graph-laplacian} to show $\bL\mathbf{1}=\bzero$.
    \item Derive $\bx^T\bL\bx=\sum_{(i,j)\in E}w_{ij}(x_i-x_j)^2$ for edges counted once.
    \item Conclude that $\bL$ is positive semidefinite for an undirected graph.
    \item Show that if a graph has two connected components, then there are two independent vectors in $\mathcal{N}(\bL)$.
\end{enumerate}
}
\label{ex:graph-laplacian-spectrum}

\addDef{Normalized Laplacians}{%
\begin{itemize}
    \item \textbf{Symmetric:}\enspace $\bL_{\text{sym}} = \bD^{-1/2} \bL \bD^{-1/2}$. (Eigenvalues in $[0, 2]$).
    \item \textbf{Random Walk:}\enspace $\bL_{\text{rw}} = \bI-\bA\bD^{-1}$ under the column-stochastic convention.
\end{itemize}
}
\label{def:normalized-laplacians}

\addRemark{%
\textbf{(Degree normalization)}\enspace The unnormalized Laplacian can let high-degree vertices dominate cuts. Normalized Laplacians measure variation relative to vertex degree.
}

\subsubsection{Spectral Clustering}

\addThm{Spectral Relaxation}{%
Minimizing a graph cut (RatioCut) is NP-hard. Spectral methods relax this to an eigenvector problem:
\begin{enumerate}
    \item Compute $k$ smallest eigenvectors of $\bL$ (starting from $\bv_2$).
    \item Embed vertices in $\fR^k$ using these eigenvectors.
    \item Cluster embedding via $k$-means.
\end{enumerate}
}

\subsubsection{Random Walks and PageRank}

\addDef{Random Walk Matrix}{%
Transition matrix $\bP = \bA\bD^{-1}$ under the column-stochastic convention.
\begin{itemize}
    \item $P_{ij} = w_{ij}/d_j$ is the probability of moving from $j$ to $i$.
    \item \textbf{Stationary Distribution:}\enspace $\boldsymbol{\pi}$ such that $\bP\boldsymbol{\pi} = \boldsymbol{\pi}$. For undirected graphs, $\pi_i = d_i / \sum d_j$.
\end{itemize}
}
\label{def:graph-random-walk}

\addExcr{%
\begin{enumerate}
    \item Show that the columns of $\bP=\bA\bD^{-1}$ sum to one.
    \item For the path graph $1-2-3$, compute $\bP$.
    \item Verify that $\pi_i=d_i/\sum_j d_j$ satisfies $\bP\boldsymbol{\pi}=\boldsymbol{\pi}$.
\end{enumerate}
}
\label{ex:graph-random-walk}

\addDef{PageRank and Google Matrix}{%
Models a random surfer on a directed web graph with teleportation:
\begin{align*}
    \mathbf{G} = d \bS + \frac{1-d}{N} \mathbf{J},
\end{align*}
\begin{itemize}
    \item $d$ is the damping factor (standard: 0.85).
    \item $\bS$ is column-stochastic and handles dangling nodes.
    \item $\mathbf{J}$ is the all-ones matrix (teleportation).
\end{itemize}
\textbf{Computation:}\enspace PageRank vector $\mathbf{r}$ is the dominant eigenvector of $\mathbf{G}$, computed via power iteration.
}
\label{def:pagerank-google}

\addRemark{%
\textbf{(Dangling nodes)}\enspace A page with no outgoing links gives a zero column before normalization. The standard fix replaces that column by the uniform probability vector before forming $\bS$.
}

\subsubsection{Exercises}

\addExcr{%
\begin{enumerate}
    \item Construct $\bL$ for a 3-vertex path graph. Find its eigenvalues by hand.
    \item Show that $\bL \mathbf{1} = \bzero$ for any graph.
    \item Prove $\bL$ is positive semidefinite using the quadratic form.
    \item Implement PageRank power iteration for a 10-node random graph.
\end{enumerate}
}

\end{document}
